% Notes / Rough CS Paper Template
% --------------------------------
% Use this as a working draft for theory/CS notes and papers.

\documentclass[11pt]{article}

% ---------- Encoding & Language ----------
\usepackage[utf8]{inputenc}
\usepackage[T1]{fontenc}
\usepackage[english]{babel}

% ---------- Page Layout ----------
\usepackage[margin=1in]{geometry}
\usepackage{setspace}
% Uncomment for double-spacing when writing a draft:
% \doublespacing

% ---------- Fonts & Math ----------
\usepackage{lmodern}
\usepackage{amsmath,amssymb,amsthm}
\usepackage{mathtools}

% ---------- Figures, Tables, Colors ----------
\usepackage{graphicx}
\usepackage{subcaption}
\usepackage{xcolor}
\usepackage{booktabs}
\usepackage{enumitem}

% ---------- Hyperlinks ----------
\usepackage{hyperref}
\hypersetup{
  colorlinks=true,
  linkcolor=blue,
  citecolor=blue,
  urlcolor=blue
}

% ---------- Bibliography (BibLaTeX stub) ----------
% If you use BibTeX instead, adjust accordingly.
% \usepackage[backend=biber,style=alphabetic,maxbibnames=99]{biblatex}
% \addbibresource{refs.bib}

% ---------- Theorem Environments ----------
\theoremstyle{plain}
\newtheorem{theorem}{Theorem}[section]
\newtheorem{lemma}[theorem]{Lemma}
\newtheorem{proposition}[theorem]{Proposition}
\newtheorem{corollary}[theorem]{Corollary}

\theoremstyle{definition}
\newtheorem{definition}[theorem]{Definition}
\newtheorem{example}[theorem]{Example}
\newtheorem{assumption}[theorem]{Assumption}

\theoremstyle{remark}
\newtheorem{remark}[theorem]{Remark}
\newtheorem{noteenv}[theorem]{Note}

% ---------- Custom Commands for CS/Math ----------
\newcommand{\R}{\mathbb{R}}
\newcommand{\N}{\mathbb{N}}
\newcommand{\Z}{\mathbb{Z}}
\newcommand{\E}{\mathbb{E}}
\newcommand{\Var}{\operatorname{Var}}
\newcommand{\Prob}{\mathbb{P}}

\newcommand{\bigO}{\mathcal{O}}
\newcommand{\smallo}{o}

\newcommand{\1}{\mathbf{1}}
\newcommand{\vect}[1]{\mathbf{#1}}
\newcommand{\abs}[1]{\left\lvert #1 \right\rvert}
\newcommand{\norm}[1]{\left\lVert #1 \right\rVert}
\newcommand{\set}[1]{\left\{ #1 \right\}}
\newcommand{\paren}[1]{\left( #1 \right)}
\newcommand{\bracket}[1]{\left[ #1 \right]}

% ---------- TODOs / Draft Notes ----------
\newcommand{\TODO}[1]{\textcolor{red}{[TODO: #1]}}
\newcommand{\idea}[1]{\textcolor{purple}{[Idea: #1]}}
\newcommand{\fixme}[1]{\textcolor{orange}{[Fix: #1]}}

\setlength{\parindent}{0pt}

% ---------- Title Information ----------
\title{Working Notes / Rough Draft\\[4pt]\large (Self-supervised CoT Monitoring for Post Training)}
\author{Dhruman Gupta, Aritra Das, Debayan Gupta}
\date{\today}

\begin{document}

\maketitle

% Optional abstract for paper-style notes
\begin{abstract}
These are working notes toward a CS/theory paper. Summarize the problem, main ideas, and key results here once they stabilize.
\end{abstract}

\tableofcontents
\newpage

% ================== SECTION TEMPLATE ==================

\section{Introduction}

% High-level problem statement, motivation, and context.
% Use this section to collect informal thoughts and references.

CoT is bad, lies, opaque etc. We believe that this is due to unoptimal learning (since CoT is not rewarded, only outputs are). Moreover, this leads to "shortcuts" that circumvent instructions and cause reward hacking.

\subsection{Literature Review}

Let $\pi_\theta$ be the current policy. While post training, the standard GRPO reward is used to optimise the model:

\definition{GRPO Reward, Advantage, and Objective}\\
For each prompt we sample a \emph{group} of $G$ trajectories with scalar rewards $\{R_i\}_{i=1}^G$ (e.g., RM score plus any KL penalty). GRPO first defines a group-normalized reward (used as an advantage) for trajectory $i$:
$$
\tilde R_i \;=\; \frac{R_i - \bar R}{\sigma_R + \epsilon}, \qquad
\bar R = \frac{1}{G}\sum_{j=1}^G R_j,\quad
\sigma_R = \sqrt{\frac{1}{G}\sum_{j=1}^G (R_j - \bar R)^2},
$$
where $\epsilon > 0$ is a small constant for numerical stability. This quantity is then broadcast to all tokens $t$ in trajectory $i$ to define the per-token group-relative advantage
$$
\hat A_{i,t} \coloneqq \tilde R_i.
$$
Let $r_{i,t}(\theta)$ be the PPO-style importance ratio at token $t$ of trajectory $i$,
$$
r_{i,t}(\theta) \;=\; \frac{\pi_\theta(o_{i,t} \mid q_i, o_{i,<t})}{\pi_{\theta_{\mathrm{old}}}(o_{i,t} \mid q_i, o_{i,<t})},
$$
and let $\varepsilon$ denote the clipping parameter and $\beta$ the KL penalty weight. A standard form of the GRPO clipped surrogate objective is then
$$
J_{\mathrm{GRPO}}(\theta)
= \mathbb{E}\!\left[
\frac{1}{G}\sum_{i=1}^G \frac{1}{|o_i|}\sum_{t=1}^{|o_i|}
\min\!\Big(
r_{i,t}(\theta)\,\hat A_{i,t},\;
\operatorname{clip}\!\big(r_{i,t}(\theta), 1-\varepsilon, 1+\varepsilon\big)\,\hat A_{i,t}
\Big)
\right]
- \beta\,D_{\mathrm{KL}}\!\big(\pi_\theta \,\|\, \pi_{\mathrm{ref}}\big),
$$
which matches PPO's clipped objective but uses the group-relative, variance-normalized advantage instead of a critic-based advantage.\\

As evident, the reward is only assigned to the output tokens, and the CoT is not rewarded. While this vanilla formulation induces biases on the CoT implicitly (for example length), this is addressed by newer work in the field (Dr. GRPO, etc).

\section{Method}

Let $\pi_{\theta_\mathrm{old}}$ be the current policy. Then, generate rollouts using $\pi_{\theta_\mathrm{old}}$ $x_1, x_2, \ldots, x_G$, where $x_i = (o_i, c_i)$ is an output-CoT pair. Now, we obtain the GRPO reward for each $o_i$.

% The score is given by a prompt as a "Verifier", which checks whether any reward hacking etc has occurred. Can use meta verifier here, same method as DeepSeek. Now, this leads to two methods:

\subsection{Method 1 - DeepSeekMath-V2 method}
\begin{enumerate}
  \item Train a verifier model $v_\phi$ to score the trajectory pair $x_i$.
  \item Train a meta-verifier model $mv_\phi$ to score the verifier model $v_\phi$.
  \item Then train the model to maximize the reward of the verifier and meta-verifier model $v_\phi$ and $mv_\phi$.
  \item Naturally the CoT will be optimized for, since it is rewarded.
  \item Add this with the accuracy loss to ensure that the model is still accurate, but is nudged towards the desired behavior.
  \item Pros and cons as the DeepSeek paper. Needs to be more thought out.
\end{enumerate}

\subsection{Method 2 - Old weight method}
\begin{enumerate}
  \item Use the model $\pi_{\theta_\mathrm{old}}$ to then review the pair $x_i$ to monitor for reward hacking, and assign score $tr_i \in \{0, 0.5, 1\}$ to each trajectory pair $x_i$.
  \item Use this reward as another reward along with the existing benchmark reward.
  \item Do GRPO.
  \item Conjecture: we are nudging for higher accuracy and CoT faithfulness. Suppose we give weight 1 to the accuracy reward and weight $\alpha$ to the CoT reward ($\alpha < 1$). If the model fails to supervise the CoT correctly (i.e the reward is flawed), then the accuracy will drop (and hence the overall reward), thus incentivizing the model to supervise the CoT correctly.

\end{enumerate}

\section{Miscellaneous Notes}

Anything that does not fit naturally elsewhere: heuristics, empirical
observations, notation debates, etc.

% ------------------------------------------------------

% \printbibliography % Uncomment if using BibLaTeX

\end{document}

